\documentclass[12pt]{article}

\usepackage[margin=1in]{geometry}
\usepackage{amsmath,amssymb,amsthm,mathtools}
\usepackage{booktabs}
\usepackage{microtype}
\usepackage{hyperref}
\usepackage[round,authoryear]{natbib}

\hypersetup{hidelinks,pdfauthor={Anonymous},pdftitle={A Correction to the Two-Type Refund Result in Ringbom and Shy (2004)}}
\newtheorem{proposition}{Proposition}
\newtheorem{corollary}{Corollary}
\newtheorem{lemma}{Lemma}
\newtheorem{remark}{Remark}

% Economics Bulletin initial-submission formatting: centered bold 14pt section headings.
\makeatletter
\renewcommand\section{\@startsection{section}{1}{\z@}%
  {-2.0ex \@plus -0.5ex \@minus -.2ex}%
  {1.0ex \@plus .2ex}%
  {\centering\normalfont\bfseries\fontsize{14}{16}\selectfont}}
\renewcommand\subsection{\@startsection{subsection}{2}{\z@}%
  {-1.5ex \@plus -0.5ex \@minus -.2ex}%
  {0.8ex \@plus .2ex}%
  {\centering\normalfont\bfseries\fontsize{14}{16}\selectfont}}
\makeatother
\newcommand{\E}{\mathbb{E}}
\newcommand{\argmax}{\operatorname*{arg\,max}}

\begin{document}

% Economics Bulletin generates the title page from portal metadata.
% The submitted PDF must begin with Section 1 and contain no abstract/title page.
\section{Introduction}

\citet{ringbomShy2004} study a seller who chooses a partial refund for advance reservations. In their two-type section, a high-show-probability type is just willing to reserve at $r_H$, while both types are willing at the larger rate $r_L$. Equation (5) correctly compares the corresponding endpoint profits. The condition printed in equation (6) and Proposition 1, however, understates the high-type loss from raising the refund: the term must be divided by $1-\sigma_L$ (Ringbom and Shy, 2004, Eq. (6), p. 3).

The omission can reverse the endpoint ranking even when both rates lie strictly inside the feasible interval. At the exact parameter point reported below, the printed condition favors $r_L$, while primitive expected profits favor $r_H$. A second issue concerns global choice. The endpoint comparison does not establish that either rate is optimal when a threshold is outside $[0,1]$ or when both feasible endpoint profits are below the zero-profit option of taking no reservations. We give a complete correspondence for the source's strategy set and tie convention.

The error is confined to the two-type comparison. We independently rederive the later uniform-type formulas and welfare results on their active domain. Those results survive, subject to their domain qualifications, but two numerical private-refund entries in Table 1 are incorrect. We report those entries so that the correction does not overstate what survives.

\section{Two-type model}

There are two consumer types, $H$ and $L$, with show probabilities $0<\sigma_L<\sigma_H<1$ and population shares $\alpha_H,\alpha_L>0$, where $\alpha_H+\alpha_L=1$. The seller chooses a refund rate $r\in[0,1]$. A consumer who shows receives value $\beta$ and pays price $p$; a consumer who cancels receives refund $rp$. The seller incurs service cost $c$ and obtains salvage value $s$ after cancellation. Consumers reserve when indifferent, as in the weak participation condition used in the source's continuum section.

Type $i$'s expected utility is
\begin{equation}
 U_i(r)=\sigma_i(\beta-p)-(1-\sigma_i)(1-r)p
       =\sigma_i(\beta-rp)-p(1-r).
 \label{eq:utility}
\end{equation}
For $p>0$, this is strictly increasing in $r$. The indifference rate and expected seller profit per reservation are
\begin{align}
 r_i&=\frac{p-\beta\sigma_i}{p(1-\sigma_i)},
 \label{eq:threshold}\\
 g_i(r)&=\sigma_i(p-c)+(1-\sigma_i)((1-r)p+s-c)\\
       &=p-c-(1-\sigma_i)(rp-s).
 \label{eq:unitprofit}
\end{align}
At a type's indifference rate,
\begin{equation}
 g_i(r_i)=q_i:=\sigma_i(\beta-s)+s-c.
 \label{eq:q}
\end{equation}
Moreover,
\begin{equation}
 r_L-r_H=\frac{(\beta-p)(\sigma_H-\sigma_L)}{p(1-\sigma_H)(1-\sigma_L)}.
 \label{eq:thresholdgap}
\end{equation}
Thus the ordering $r_H<r_L$ requires $\beta>p$. Thresholds must still be checked against the strategy set; a raw negative threshold is not a feasible refund.

\section{Correct endpoint comparison}

Suppose $0\leq r_H<r_L\leq1$. Under weak participation, only type $H$ reserves at $r_H$, and both types reserve at $r_L$. Let $\Pi_H=n\alpha_Hg_H(r_H)$ and $\Pi_{HL}=n[\alpha_Hg_H(r_L)+\alpha_Lg_L(r_L)]$, where $n$ is the potential market size.

\begin{proposition}[Corrected endpoint comparison]
\label{prop:corrected}
For feasible ordered thresholds $0\leq r_H<r_L\leq1$,
\begin{equation}
 \frac{\Pi_{HL}-\Pi_H}{n}
 =\alpha_L[\sigma_L(\beta-s)-c+s]
 -\frac{\alpha_H(\sigma_H-\sigma_L)(\beta-p)}{1-\sigma_L}.
 \label{eq:correctgap}
\end{equation}
Consequently, $r_L$ weakly dominates $r_H$ among these two endpoint offers if and only if the right-hand side of \eqref{eq:correctgap} is nonnegative. Equality holds exactly when the two endpoint profits tie.
\end{proposition}

\begin{proof}
At $r_L$, type $L$ earns $q_L$ by \eqref{eq:q}. Type $H$'s change in unit profit between the rates is $-(1-\sigma_H)p(r_L-r_H)$. Substituting \eqref{eq:thresholdgap} gives
\[
 (1-\sigma_H)p(r_L-r_H)
 =\frac{(\sigma_H-\sigma_L)(\beta-p)}{1-\sigma_L}.
\]
Multiplying by the high-type mass and adding the low-type contribution yields \eqref{eq:correctgap}. The final equivalence follows from the sign of the exact profit difference.
\end{proof}

Equation (6) in the source omits $1/(1-\sigma_L)$. The difference is not a normalization: it is the high-type cost of moving from $r_H$ to $r_L$ after substituting the distance between the two indifference rates.

\paragraph{Exact counterexample.}
Set
\[
 (\beta,p,c,s)=\left(1,\frac35,\frac1{50},0\right),\quad
 (\alpha_H,\alpha_L)=\left(\frac45,\frac15\right),\quad
 (\sigma_H,\sigma_L)=\left(\frac{11}{20},\frac7{20}\right).
\]
These parameters satisfy $\beta>p>c\geq s\geq0$. The thresholds are $r_H=5/27$ and $r_L=25/39$, so both are feasible and strictly ordered. Direct primitive-payoff evaluation gives
\[
 \frac{\Pi_H}{n}=\frac{53}{125},\qquad
 \frac{\Pi_{HL}}{n}=\frac{509}{1300},\qquad
 \frac{\Pi_{HL}-\Pi_H}{n}=-\frac{211}{6500}.
\]
The printed condition has left side minus right side $1/500>0$ and therefore selects $r_L$, although $r_H$ earns more. All values are exact rationals; a separate checker computes expected seller cash flows directly from the two states.

\section{Global refund choice and the first corollary}

For any refund, define the primitive profit schedule
\begin{equation}
 \Pi(r)=n\sum_{i\in\{H,L\}}\alpha_i\mathbf{1}\{U_i(r)\geq0\}g_i(r),\qquad r\in[0,1].
 \label{eq:profit}
\end{equation}
The indicator in \eqref{eq:profit} makes explicit the weak participation convention. On any interval with a fixed, nonempty participating set $A$,
\begin{equation}
 \Pi'(r)=-np\sum_{i\in A}\alpha_i(1-\sigma_i)<0.
 \label{eq:slope}
\end{equation}
Thus profit falls strictly between threshold events. The expressions displayed in the source's equation (4) substitute threshold values; they are endpoint profits rather than the profit at every interior refund in each region.

\begin{proposition}[Complete feasible refund correspondence]
\label{prop:global}
Let $\beta,p>0$, $c,s\geq0$, $0\leq\sigma_L<\sigma_H<1$, $\alpha_i\geq0$ with $\alpha_H+\alpha_L>0$, and use \eqref{eq:utility}--\eqref{eq:profit}. The weights need not be normalized; normalizing them only rescales profit. Let $J=\{i:\alpha_i>0\}$, $\rho_i=(p-\beta\sigma_i)/(p(1-\sigma_i))$, and
\[
 \mathcal C=\{0\}\cup\{\rho_i:i\in J,\ 0\leq\rho_i\leq1\},\quad
 M=\max_{r\in\mathcal C}\Pi(r),\quad
 \mathcal C^*=\{r\in\mathcal C:\Pi(r)=M\}.
\]
Define $\rho_0=\min_{i\in J}\rho_i$. When all positive-mass thresholds are positive and $M=0$, the argmax is
\[
 \mathcal C^*\cup\{r\in[0,1]:r<\rho_0\}.
\]
In every other case the argmax is exactly $\mathcal C^*$.
\end{proposition}

\begin{proof}
Each utility is affine and strictly increasing in $r$, so a type joins at its raw threshold if that threshold lies in $[0,1]$, is already present at $r=0$ if its threshold is negative, and never joins if its threshold exceeds one. The candidate set includes zero and every in-domain participation threshold of a positive-mass type. Between consecutive candidates, either nobody participates and profit is zero, or a fixed positive mass participates and \eqref{eq:slope} makes profit strictly decreasing. Hence a positive-profit maximum can occur only at a candidate. If every positive-mass threshold is positive, the initial no-reservation set $\{r\in[0,1]:r<\rho_0\}$ has profit zero; it belongs to the argmax exactly when $M=0$. Weak acceptance evaluates threshold equalities at the threshold itself, including $r_i=0$ and $r_i=1$. In particular, at $\beta=p$ both thresholds equal one; for $\beta<p$, positive-show-probability types have thresholds above one. The argument uses primitive utility and does not divide by $\beta-rp$, so it also covers the boundary $\beta-rp=0$.
\end{proof}

When $0<r_H<r_L<1$ and both masses are positive, the candidate set is $\{0,r_H,r_L\}$. The following implication corrects the global reading of the source's Corollary 1.

\begin{corollary}[Both types served]
On $0<r_H<r_L<1$, the seller selects $r_L$ if and only if
\begin{equation}
 \Pi_{HL}\geq \max\{0,\Pi_H\}.
 \label{eq:correctioncor}
\end{equation}
At equality, the optimal-refund correspondence includes both tied candidates (and, when both endpoints and no booking earn zero, the no-reservation interval).
\end{corollary}

The no-booking option can matter even near the paper's high-price region. For $\beta=1$, $p=999/1000$, $c=9/10$, $s=0$, equal masses, and $(\sigma_H,\sigma_L)=(1/2,1/10)$, the thresholds are $998/999$ and $8990/8991$. Exact profits are $-1/5$ at $r_H$ and $-2701/4500$ at $r_L$, while taking no reservations yields zero. Thus feasibility and endpoint ranking alone do not establish global service of both types.

\section{What survives in the uniform model}

The later continuous-type analysis is separate from the two-type endpoint comparison. For a uniform show probability $\sigma\sim U[0,1]$, on $\beta>p>c\geq s\geq0$ the reservation cutoff is
\[
 \widehat\sigma(r)=\frac{p(1-r)}{\beta-rp},\qquad
 \frac{\pi(r)}n=\int_{\widehat\sigma(r)}^1[p-c-(1-\sigma)(rp-s)]\,d\sigma.
\]
Direct differentiation gives
\[
 \frac{\pi'(r)}n=\frac{p(\beta-p)}{(\beta-rp)^2}
 \left[g_{\widehat\sigma(r)}(r)-\frac{\beta-p}{2}\right].
\]
The bracket is strictly decreasing in $r$, establishing a unique constrained private optimum. Solving for its root gives
\[
 r_{\mathrm{int}}=\frac{\beta[3p-2(c-s)]-2ps-\beta^2}{p(\beta-2c+p)},\qquad
 \bar r=\max\{0,r_{\mathrm{int}}\}.
\]
Recomputing this formula and primitive profits gives $\bar r=8/13$ rather than the printed $0.250$ at $(\beta,p,c,s)=(1,1/2,1/10,1/10)$, and $\bar r=2/3$ rather than $0.357$ at $(1,1/2,0,0)$. Each printed rate has a positive derivative bracket, and the exact direct-profit gains at the corrected rates are $361/35280$ and $863041/86382368$ respectively.

The planner's objective is
\[
 \frac{W(r)}n=\int_{\widehat\sigma(r)}^1[(\beta-s)\sigma-(c-s)]\,d\sigma,
\]
because refund transfers cancel between consumers and the seller. Its optimal cutoff is $(c-s)/(\beta-s)$, giving
\[
 r^*=1-\frac{(c-s)(\beta-p)}{p(\beta-c)}.
\]
On the stated strict domain the source's welfare comparisons and loss formulas follow independently of equation (6). Write $\widetilde p=\beta(\beta+2c-2s)/(3\beta-2s)$. The welfare loss is
\[
 \Delta W=\begin{cases}
 \displaystyle\frac{n(\beta-p)^2}{8(\beta-s)},&p>\widetilde p,\\[6pt]
 \displaystyle\frac{n[p(\beta-s)-\beta(c-s)]^2}{2\beta^2(\beta-s)},&p\leq\widetilde p.
 \end{cases}
\]
In particular, $r^*>\bar r$ on the strict domain. The strict inequality needs $p>c$: at the excluded boundary $p=c,s=0$, both rates can equal zero. All remaining Table 1 values and all Table 2 welfare-loss entries reproduce at the published precision when evaluated at the primitive optimum. The uniform cutoff expression is not used at $\beta-rp=0$; boundary cases require direct utility rather than division by a zero denominator.

\section{Conclusion}

The correct two-type endpoint comparison adds the factor $1/(1-\sigma_L)$ to the high-type loss term in equation (6). Its omission can reverse the ranking of feasible refund endpoints. A complete global choice also evaluates threshold feasibility, the relevant tie rule, and the no-booking alternative. These corrections qualify Proposition 1 and Corollary 1, while the later uniform welfare analysis remains valid on its active domain. Two numerical private-refund entries in Table 1 require correction.

\begin{thebibliography}{9}
\bibitem[Ringbom and Shy(2004)]{ringbomShy2004}
Ringbom, S. and O. Shy (2004).
\newblock Advance booking, cancellations, and partial refunds.
\newblock \emph{Economics Bulletin}, 13(1), 1--7.
\newblock Publisher Version of Record: \url{https://www.accessecon.com/pubs/EB/2004/Volume13/EB-04M20001A.pdf}.
\end{thebibliography}

\end{document}
